\def\ChapterCode{JUS}
\chapter
[\CO{[JUS]} Rechtliche Grundlagen (.at)]
{Rechtliche Grundlagen (.at)}
\label{chp:JUS}
% \AasRsN{RS~III}

\setcounter{tocdepth}{10}
\mtcsetdepth{minitoc}{3}
\ChapterIntro
{%
Dieses Kapitel behandelt die allgemeinen und
tätigkeitsspezifischen rechtlichen Grundlagen. 

Anhand von konkreten Beispielen vermittelt es die notwendige
Kompetenz und Problemlösungsstrategien um im Alltag
auftretende rechtliche Probleme und Fragestellungen zu
lösen.
}
{%
\begin{description}
  \item[Maintainer:] Christof Koller
  \item[Autoren:] Christof Koller et. al.
%   \item[Reviewer:] --
  \item[Version:] {\DocVersion} ({\DocVersionInfo})
  \item[SHA1:] {\DocGitCommit}
\end{description}

{\TxTeilkapitelAASS}
}
% \minitoc
\mtcsetdepth{minitoc}{\value{DefaultMinitocDepth}}
% \AuthNotes{
% Changelog:
% 2009-02-20: GABS: Fertig für Kurs 2009/02
% 
% 2009-02-25: EDIT Koller
% 
% 2009-04-22: GABS: Abkürzungen \enquote{Pat.}, \enquote{KA}, \enquote{KH}
% bereinigt und gegen Langform ausgetauscht, wie besprochen.}





\clearpage
\section{Allgemeine Rechtsgrundlagen}



\subsection{Struktur der Österreichischen Rechtsordnung}



%\begin{tabular}{| p{2.2cm} | p{2.2cm} p{2.2cm} p{2.2cm} | p{2.2cm} p{2.2cm}
% |}\toprule

\TextSlide{\TxBeschreibung}{
\begin{table}[H]
\TableBaseModification
\Captionof{table}{Rechtsordnung}
\begin{tabularx}{\linewidth}[H]{X X|X|X||X|X||}

& \multicolumn{5}{c}{\EE{$\Swarrow$ EU-Recht $\Searrow$}}
\\\addlinespace

& \multicolumn{5}{c}{\EE{Struktur der Österreichischen Rechtsordnung}}
\\\addlinespace

{ }
& \multicolumn{3}{c}{\centering Öffentliches Recht}
& \multicolumn{2}{c}{\centering {Privatrecht} }
\\\hhline{~-----}

\multicolumn{1}{p{2cm}||}{Rechts\-gebiete }
& Strafrecht
& VwR
& VfR
& Zivilrecht
& Sonder\-privat\-rechte \\\hhline{~-----}

\multicolumn{1}{p{2cm}||}{Wichtigste Gesetzesbereiche} &
StGB, StPO &
Berufs\-ge\-setze,\newline KAKuG, UbG &
B-VG; StGG &
ABGB &
Arbeits\-rechts\-ge\-setze,
Unter\-nehmens\-ge\-setz\-buch\\\hhline{~-----}

\end{tabularx}
\end{table}
%\ncbarr[arrows=->]{EU}{OR}


Die österreichische Rechtsordnung ist ein
zusammenhängendes Geflecht bestehend aus zwei
Hauptbereichen, dem öffentlichen Recht und dem Privatrecht.
Während das öffentliche Recht die Beziehung zwischen Staat
und Bürger regelt, bestimmt das Privatrecht die Verhältnisse
unter den Bürgern (= Privaten) selber:

\begin{itemize}
  \item \protect\T[Recht!Öffentliches@Recht!Oeffentliches]{Öffentliches Recht}:
  Staat ${\leftrightarrow}$ Bürger
  \item \T{Privatrecht}: Bürger ${\leftrightarrow}$ Bürger
\end{itemize}

Durch den Beitritt zur \E{EU} wurde ein weiterer
Gesetzgeber geschaffen, der in beide Rechtsgebiete einwirken
kann. Unabhängig davon lässt sich jedes von beiden
Rechts\-ge\-bieten in folgende Bereiche unterteilen:
}{
\begin{itemize}
  \item Öffentliches Recht
  \begin{itemize}
    \item Strafrecht
    \item Verwaltungsrecht
    \item Verfassungsrecht
  \end{itemize}
  \item Privatrecht
  \begin{itemize}
    \item Zivilrecht
    \item Sonderprivatrechte
  \end{itemize}
\end{itemize}
}{}{}{}

\Text{Öffentliches Recht}
{%
\begin{itemize}
    \item \T{Strafrecht}: Erfasst jene Delikte, die
    der Gesetzgeber als derart schwerwiegend erachtet hat,
    dass sie nicht nur (idR.) vom Staatsanwalt vor Gericht
    zu verfolgen sind,  sondern auch (in den meisten Fällen)
    mit Haftstrafen  sanktionierbar sind. Der Staatsanwalt
    vertritt hierbei  den Staat als Ankläger. Die
    eingehobene Geldstrafe  erhält der Staat, nicht der
    Verletzte.
    
    \item \T{Verwaltungsrecht}: Umfasst Regelungen
    die sich idR. nur an bestimmte Personenkreise richten
    und von Verwaltungsbehörden (in 1. Instanz) und
    von Verwaltungsgerichten (in 2. Instanz)
    bei Verstoß insofern sanktioniert werden, dass
    (vorrangig) Geldstrafen verhängt werden.
    
    \item \T{Verfassungsrecht}: Jener Teil der
    Rechtsordnung, der die grundlegenden Bestimmungen über
    Staatsgewalt, ihre Organe und ihre Verhältnisse zu den
    Staatsbürgern sowie Grundsätze der Rechtsordnung
    enthält. Er umfasst weiters Grundprinzipien und
    Grundrechte einer Rechtsordnung.
    
\end{itemize}
}{}{}{}{}



\Text{Privatrecht}
{%
\begin{itemize}
    \item \T{Zivilrecht}: Enthält jene allgemeinen Normen,
    die die Rechtsverhältnisse zwischen gleichgestellten
    Bürgern regeln.
    Bei Verstößen muss der Verletzte selber vor Gericht sein
    Recht einfordern (kein Staatsanwalt), wodurch er vom
    Schädiger Ersatzleistungen für den Schaden fordern kann
    (keine Strafe).
    
    \item \T{Sonderprivatrechte}: eigenständige
    Rechtsgebiete, die mittels selbstständigen Normen
    besondere Rechtsverhältnisse zwischen den Bürgern
    regeln. Zu den wichtigsten zählen das Arbeits- und das
    Unternehmensrecht.
    
\end{itemize}
}{}{}{}{}

\Text{Zugang zu den Gesetzestexten}
{Die Gesetzestexte können im Internet vom Rechtsinformationssystem des Bundes
(RIS) unter der URL \url{http://www.ris.bka.gv.at/} abgerufen werden.

\noindent\includegraphics[width=\linewidth]{./images/JUS/RisHeader.jpg}
}{}{}{}{}

\subsection{Vertragsverhältnis im Rettungswesen}
\label{topic:vetrag-rettungswesen}



\TextSlide{\TxBeschreibung}{%
\begin{figure}[H]
%     \AuthNotes{KOLC: Graphik gehört rotiert.
%
%     GABS: ok, dauert aber noch ein bisserl. Wenn Inhalt
%     passt bitte etwas Geduld.}
\begin{center}
\includegraphics[width=\linewidth]{./images/gabriel-sebastian-aass/recht_vertrag.pdf}
\end{center}
\Captionof{figure}{Vertragsverhältnisse im Rettungsdienst}
\end{figure}
Das Rechtsverhältnis im Rettungswesen besteht
grundsätzlich aus drei Parteien: dem Patienten, der
Rettungsorganisation (Einrichtung) und dem (Not-) Arzt oder
Sanitäter (Personal).
(In einer Krankenanstalt ist das Verhältnis ähnlich, mit dem
Unterschied, dass an die Stelle der Rettungsorganisation die
Krankenanstalt und an jene des Sanitäters der Krankenpfleger
tritt.)

Als wichtigstes Verhältnis in dieser
Dreiecksbeziehung ist das \EE{Verhältnis zwischen Patient
und Rettungsorganisation} zu bezeichnen, da direkter
(Vertrags-)Partner zum Patient niemals der Sanitäter,
sondern stets die von ihm vertretene Organisation ist. Man
nimmt nämlich an, dass ein Patient, wenn er eine
\enquote{Rettung} verständigt, nicht mit dem einzelnen
Einsatzpersonal in rechtlichen Kontakt treten möchte,
sondern mit der Gesamtheit einer Organisation. (Er bestellt
ja nicht dieses eine Fahrzeug mit jenem speziellen
Personal.) Das Rechtsverhältnis kann nun auf zwei
verschiedene Arten zustande kommen:

\begin{itemize}
  \item
  \T[Heilbehandlung!-svertrag]{Heilbehandlungsvertrag}: Der
  Patient selber oder sein
 zuständiger Vertreter wünscht behandelt zu werden
  \item \EE{\T{Geschäftsführung ohne Auftrag} im Notfall}:
  der Patient befindet sich in einem derartigen Notfall,
  dass er sich nicht mehr äußern kann und dringend eine
  Behandlung benötigt (Bsp.:
  Bewusstlosigkeit)
\end{itemize}

Aus beiden Verhältnissen schuldet die
Rettungsorganisation die höchstmögliche Sorgfalt, aber
niemals einen bestimmten Behandlungserfolg. Weiters hat sie
vorrangig nach dem Willen des Patienten vorzugehen und wenn
dieser fehlt, nach seinem hypothetischen Willen.

Im \EE{Verhältnis zwischen Patient und Personal}
fungiert das Personal bloß als Gehilfe zur Erfüllung der
Pflichten seiner Organisation gegenüber dem Patienten,
wodurch man ihn auch als
\T[Erfüllungsgehilfe]{Erfüllungsgehilfen} bezeichnet. Er ist
nicht Vertragspartner des Patienten, muss aber mit der
gleichen Sorgfalt, wie es von der Organisation verlangt
wird, arbeiten. Kommt es zu Schädigungen des Patienten durch
Handlungen des Personals, kann der Geschädigte deliktisch
gegen den Erfüllungsgehilfen und vertraglich gegen die
Organisation vorgehen. (I.\,d.\,R. wird der Patient die
Organisation verklagen, da ihm hierbei aufgrund des
Vertrages mehrere gesetzliche Erleichterungen zustehen.)

Das \EE{Verhältnis zwischen Rettungsorganisation und
Personal} lässt sich auf drei verschiedene Arten begründen:

\begin{itemize}
    \item \E{Arbeitsvertrag}: Der (Not-)Arzt/Sanitäter ist
    bei der Organisation als Hauptamtlicher berufsmäßig
    beschäftigt.
    
    \item \E{Vereinsmitgliedschaft}: Der
    (Not-)Arzt/Sanitäter ist bei der Organisation als
    Vereinsmitglied und somit als
    Ehrenamtlicher/Freiwilliger tätig.
    
    \item \E{Zivildienst}: Der Sanitäter leistet seinen
    Zivildienst bei der Organisation.
    
\end{itemize}

Aus diesem Verhältnis heraus müssen manche der Pflichten,
welche die Organisation gegenüber dem Patienten treffen,
auch durch das Personal erfüllt werden. Gleichzeitig treffen
auch die Organisation und das Personal Rechte und Pflicht im
Umgang miteinander.
}{
\begin{itemize}
  \item Verhältnis zwischen Patient und Rettungsorgansiation
  \begin{itemize}
    \item Heilbehandlungsvertrag
    \item Geschäftsführung ohne Auftrag im Notfall
  \end{itemize}
  \item Verhältnis zwischen Patient und Personal
  \begin{itemize}
    \item Personal ist Erfüllungsgehilfe
  \end{itemize}
  \item Verhältnis zwischen Rettungsorganisation und
  Personal
  \begin{itemize}
    \item Arbeitsvertrag
    \item Vereinsmitgliedschaft
    \item Zivildienst
  \end{itemize}
\end{itemize}
}{}{}{}

\subsection{Straf- und Haftbarkeitsvoraussetzungen}
\label{topic:strafbarkeitspruefung}
\label{topic:haftbarkeitspruefung}
\label{topic:strafbarkeitsvoraussetzung}
\label{topic:haftbarkeitsvoraussetzung}



\TextSlide{Voraussetzungen}
{Damit es zur Strafbarkeit
(Strafrecht) und/oder zur Haftbarkeit (Zivilrecht, Schadenersatz) für eine
gesetzte Handlung kommt, müssen die folgenden
überblicksartig dargestellten \E{Voraussetzungen}
vollständig erfüllt sein:

\begin{enumerate}
  \item \EE{Tatbestandserfüllung} bzw. \EE{Schaden}:
  Erfüllung der rechtlichen Beschreibung einer
  Lebenssituation auf die eine Rechtsfolge/Strafe angeordnet
  wird.
  \item \EE{Kausalität}: Die gesetzte Handlung war der Grund
  für den Eintritt der Tatbestandserfüllung
  \item \EE{Rechtswidrigkeit}: Die gesetzte Handlung war
  objektiv falsch ${\rightarrow}$ Urteil über Tat
  ${\rightarrow}$ Verstoß gegen gesetzliche Verbote bzw.
  Gebote, vertraglichen Bestimmungen oder die guten Sitten
  
  \EE{Rechtfertigungsgründe}:  Der begangenen Tat kann kein
  Urteil gemacht werden, da sich ein maßgerechter Mensch in 
  derselben Situation gleich verhalten
  hätte:
  \begin{compactitem}
    \item Heilbehandlung
    \item Einwilligung
    \item Notwehr
    \item rechtfertigender Notstand
    \item Anhalterecht Privater
  \end{compactitem}
  \item \EE{Verschulden}: Die gesetzte Handlung war
  subjektiv falsch ${\rightarrow}$ Urteil über
  Täter
  
  \EE{Entschuldigungsgründe}: Über den Täter kann kein
  Urteil gemacht werden, da er sich nach seinen persönlichen
  Fähigkeiten und Kenntnissen in dieser Situation nicht
  anders verhalten hätte können:
    \begin{compactitem}
      \item Notwehrüberschreitung im Affekt
      \item entschuldigender Notstand
      \item (Unzurechenbarkeit)
    \end{compactitem}
  
  Die einzelnen Rechtfertigungs- und Entschuldigungsgründe
  werden separat behandelt.
  
\end{enumerate}

Mangelt es auch nur an
einer von diesen Voraussetzungen, kommt es zu keiner
Straf-/ bzw. Haftbarkeit. Liegen die Voraussetzungen
hingegen vor, kommt es im Strafrecht zu einer \I{Strafe}
(Haft- oder Geldstrafe zu Gunsten des \E{Staates}), bzw. 
im Zivilrecht zu einem \I{Schadenersatz} (z.\,B.
Heilungskosten, Verdienstentgang, Schmerzengeld, Ersatz für
entgangenen Gewinn oder realen/ideellen Schaden).

}{
\begin{enumerate}
  \item \EE{Tatbestandserfüllung} bzw. \EE{Schaden}
  \item \EE{Kausalität}
  \item \EE{Rechtswidrigkeit}

  Rechtfertigungsgründe:
  \begin{compactitem}
    \item Heilbehandlung
    \item Einwilligung
    \item Notwehr
    \item rechtfertigender Notstand
    \item Anhalterecht Privater
  \end{compactitem}
  \item \EE{Verschulden}
  
  {Entschuldigungsgründe}:
    \begin{compactitem}
      \item Notwehrüberschreitung im Affekt
      \item entschuldigender Notstand
      \item (Unzurechenbarkeit)
    \end{compactitem}  
\end{enumerate}
}



\Text{Fallbeispiele}
{
\begin{itemize}
  \item Fall 13: Sachbeschädigung
  \dotfill\FullPrettyRef{topic:jus-langer-dienst-fall-13}
  \item Fall 16: Strafbare Handlungen
  \dotfill\FullPrettyRef{topic:jus-langer-dienst-fall-16}
\end{itemize}
}
{}
{}
{}
{}


\clearpage

%%%%%%%%%%%%%%%%%%%%%%%%%%%%%%%%%%%%%%%%%%%%%%%%%%%%%%%%%%%
%%%%%%%%%%%%%%%%%%%%%%%%%%%%%%%%%%%%%%%%%%%%%%%%%%%%%%%%%%%
%%%%%%%%%%%%%%%%%%%%%%%%%%%%%%%%%%%%%%%%%%%%%%%%%%%%%%%%%%%
%%%%%%%%%%%%%%%%%%%%%%%%%%%%%%%%%%%%%%%%%%%%%%%%%%%%%%%%%%%
%%%%%%%%%%%%%%%%%%%%%%%%%%%%%%%%%%%%%%%%%%%%%%%%%%%%%%%%%%%
%%%%%%%%%%%%%%%%%%%%%%%%%%%%%%%%%%%%%%%%%%%%%%%%%%%%%%%%%%%
\section{Besondere Rechtsgrundlagen}

%%%%%%%%%%%%%%%%%%%%%%%%%%%%%%%%%%%%%%%%%%%%%%%%%%%%%%%%%%%
%%%%%%%%%%%%%%%%%%%%%%%%%%%%%%%%%%%%%%%%%%%%%%%%%%%%%%%%%%%
%%%%%%%%%%%%%%%%%%%%%%%%%%%%%%%%%%%%%%%%%%%%%%%%%%%%%%%%%%%
%%%%%%%%%%%%%%%%%%%%%%%%%%%%%%%%%%%%%%%%%%%%%%%%%%%%%%%%%%%
%%%%%%%%%%%%%%%%%%%%%%%%%%%%%%%%%%%%%%%%%%%%%%%%%%%%%%%%%%%
\subsection{Spezielles Sanitätsrecht}



\subsubsection{Einsichts- und Urteilsfähigkeit}

\label{topic:urteilsfaehigkeit}

Die Einsichts- und Urteilsfähigkeit ist grundsätzlich bei
jedem volljährigen Menschen anzunehmen. Nur wenn
außerordentliche Situationen vorliegen, kann man sie nicht
annehmen. Solche Situationen wären:

\begin{itemize}
  \item Psychiatrische Erkrankung, vgl. auch
  \enquote{Unterbringung},
  \FullPrettyRef{topic:unterbringung-jus}
  \item Starke Berauschung (bei Alkohol ab ca.
  2,5{\textperthousand})
  \item Bewusstlosigkeit
  \item Schockzustand\footnote{Aus der juristischen
  Literatur geht nicht eindeutig hervor, ob ein \enquote{Schock}
  im Sinne psychischen Reaktion, oder der \enquote{medizinische
  Schock} im Sinne einer
  lebensbedrohlichen Kreislaufstörung (\ref{topic:schock})
  gemeint ist.}
\end{itemize}

Bei minderjährigen Patienten kommt es vorrangig auf die
Einschätzung durch das Einsatzpersonal an. 

Ist sich das Einsatzpersonal nicht sicher, ob Einsichts- und
Urteilsfähigkeit vorliegt, so ist diese bei einem mündig
Minderjährigen anzunehmen, bei einem unmündig Minderjährigen
jedoch nicht.
Nur bei schwerwiegenden Angelegenheiten muss ein Elternteil
zustimmen, selbst wenn das Kind einsichts- und urteilsfähig
wäre.
Stimmt ein Elternteil zu, der andere verweigert aber, so ist
entweder jener Teil entscheidungsbefugt der die größere
Rolle in der Erziehung des Patienten spielt (im Zweifel
die Mutter) oder das
\index{Pflegschaftsgericht}Pflegschaftsgericht ist
anzurufen. Lehnen beide ab, obwohl es sich um eine
notwendige Heilbehandlung handelt, muss das
Pflegschaftsgericht angerufen werden


\Text{Fallbeispiele}
{
\begin{itemize}
  \item Fall 4: 
  \dotfill\FullPrettyRef{topic:jus-langer-dienst-fall-4}
\end{itemize}
}
{}
{}
{}
{}



\subsubsection{Aufklärung -- Aufklärungspflicht}
\label{topic:aufklaerung}
\index{Aufklärung}
\index{Aufklärungspflicht}
Jeder Patient ist vor Setzung der Maßnahme über Umfang und
Tragweite der Heilbehandlung aufzuklären, wenn er einsichts-
und urteilsfähig ist. Die Aufklärung hat in jenem Rahmen zu
erfolgen, der einer
maßgerechten Figur unter Berücksichtigung der
individuellen Besonderheiten des Patienten und der
Situation gerecht wird. Bei der Aufklärung sind
insbesonders folgende Punkte zu berücksichtigen:

\begin{itemize}
    \item Umfang und Tragweite der  Heilbehandlung
    \item Folgen, Gefahren und Nebenwirkungen der
    Heilbehandlung
    \item Umstände die für verständigen Patienten ins
    Gewicht fallen würden
    \item umgekehrte Proportion zur Dringlichkeit der
    Heilbehandlung
  \end{itemize}

Unterbleibt die Aufklärung, wird angenommen, dass dem
Patient die Tragweite der Maßnahmen nicht bewusst ist und er
daher \E{nicht} in die Heilbehandlung einwilligt.

Seit dem 1. Dezember 2009 wurde das \EE{Recht auf
Aufklärung} durch das Inkrafttreten des Artikel~3 Abs.\,2
lit.\,a der \IX{Grundrechtecharta} der \Abk{EU} zu einem
eigenen \EE{Grundrecht}.

Unabhängig von der juristischen Bedeutung der Aufklärung ist
diese ein wesentlicher Ausdruck des Respekts gegenüber dem
Patienten und Voraussetzung für die Bildung einer
Vertrauensbasis und Führung eines professionellen Gesprächs
im Rahmen der Versorgung. 

\subsubsection{Der Revers}

\index{Revers}
\label{topic:revers}


Der Revers ist eine schriftliche Dokumentierung über die
Aussage des Patienten, eine vom Einsatz\-personal
empfohlenen Maßnahme abzulehnen. Die Verweigerung kann auch
vom zuständigen Vertreter des Patienten abgegeben werden.

Der Revers ist immer dann notwendig, wenn das
Einsatzpersonal eine Maßnahme für erforderlich erachtet, der
Patient  (oder sein zuständiger Vertreter) diese Maßnahme
aber verweigert. Er dient zur Einhaltung der gesetzlichen
Dokumentationspflicht und zur Abwendung von  klagsweise
geltendmachbaren Ansprüchen des Patienten.

\paragraph{Voraussetzungen}

\begin{itemize}
  \item Auszufüllen nach umfassender Aufklärung des
  Patienten; die Aufklärung an sich ist ebenso zu
  dokumentieren!
  \item Schriftlichkeit (Transportschein, Reversschein oder
  beigelegtes Dokument)
  \item Unterschrift des Patienten (oder seines zuständigen
  Vertreters)
  \item Nur unter vorhergehender Beiziehung eines 
  Arztes
\end{itemize}

Verweigert der Patient die Unterschriftsleistung, so ist die
Verweigerung (der Maßnahme und der Unterschriftsleistung)
zu dokumentieren und am besten durch Zeugen (weitere
Mitglieder des Einsatzpersonals) mittels
Unterschrift zu bestätigen.



\subsubsection{Selbsthilferecht}
\label{topic:selbsthilferechte}
\label{topic:notstand}
\label{topic:notwehr}
\label{topic:anhalterecht-privater}

Die Selbsthilferechte erfassen jedermann zustehende
Abwehrmaßnahmen gegen entweder vom Menschen gesetzte
Angriffe (\I{Notwehr}) oder
von menschlichen Angriffen unabhängige Gefahrensituationen
(\I{Notstand}). Diese Maßnahmen erzeugen 
Schäden im Ausmaß ähnlich des Angriffs oder der Gefahr,
gelten aber als gerechtfertigt oder
entschuldigt, je nachdem ob der Angriff oder
die Gefahr entweder tatsächlich
(\EE{objektiv}) vorliegt oder nur \EE{subjektiv}
angenommen wird\footnote{Unter \Speech{subjektiv annehmen}
versteht man eine Situation, bei der der \enquote{Verteidiger}
glaubt eine Gefahr zu erkennen, obwohl diese bei späterer
Betrachtung gar nicht bestanden hat, was aber zum
damaligen Zeitpunkt nicht erkennbar war.}. 


Unabhängig von Angriff oder Gefahr ist es möglich, bei
Vorliegen einer gerichtlich strafbaren Handlung, den Täter
durch Privatpersonen festzuhalten (\I{Anhalterecht} durch
Private).

Die Voraussetzungen für Notwehr, Notstand und Anwendung des
Anhalterechts sind in den Tafeln
\prettyref{tab:notwehr-notwehrueberschreitung} (Notwehr und
Notwehrüberschreitung), \prettyref{tab:notstand} (Notstand)
und \prettyref{tab:anhalterecht-privater} (Anhalterecht)
aufgelistet und gegenübergestellt. 



\begin{table}[H]\TableBaseModification\begin{WidePar}
\Captionof{table}{\EE{Notwehr} und
\EE{Notwehrüberschreitung} im Affekt\label{tab:notwehr-notwehrueberschreitung}}

\begin{minipage}{\linewidth}
\begin{tabularx}{\linewidth}{XX}\toprule\addlinespace
\noindent\II{Notwehr}
&
\noindent\II[Notwehrüberschreitung]{Notwehrüberschreitung im
Affekt} \\\addlinespace\midrule\addlinespace
% \begin{minipage}{\linewidth}
\begin{itemize}
  \item Gerichtliches Strafrecht (StGB\ZFootnote{{\S}\,3
  Abs. 1 StGB}) und bürgerliches Recht
  (ABGB\ZFootnote{{\S}\,1306a ABGB})
  \item Anzuwenden, wenn ein menschlicher Angriff 
  tatsächlich (\EE{objektiv}) vorliegt
\end{itemize}
% \end{minipage}
&
% \begin{minipage}{\linewidth}
\begin{itemize}
  \item Gerichtliches Strafrecht (StGB\ZFootnote{{\S}\,3
  Abs. 2 StGB}) und bürgerliches Recht
  (ABGB\ZFootnote{{\S}\,1306a ABGB})
  \item Anzuwenden, wenn ein menschlicher Angriff
  \EE{subjektiv angenommen} wird
\end{itemize}
% \end{minipage}

\\\addlinespace\midrule\addlinespace
% \begin{minipage}[t]{\linewidth}
\EE{Notwehrsituation:}

\begin{itemize}
  \item gegenwärtig oder unmittelbar drohender Angriff,
  \item rechtswidrig,
  \item von Menschen gesetzt,
  \item auf ein notwehrfähiges Gut (= Leben, Gesundheit,
  Freiheit oder
  Vermögen)
\end{itemize}
% \end{minipage}
&
% \begin{minipage}[t]{\linewidth}
\EE{Notwehrsituation:}

\begin{itemize}
  \item gegenwärtig oder unmittelbar drohender Angriff,
    \item rechtswidrig,
    \item von Menschen gesetzt,
    \item auf ein notwehrfähiges Gut (= Leben, Gesundheit,
    Freiheit oder
    Vermögen)
    \item \EE{Furcht} oder Angst der in
    Notwehrüberschreitung handelnden Person
\end{itemize}
% \end{minipage}

\\\addlinespace
% \begin{minipage}{\linewidth}
\EE{Notwehrhandlung}:
\begin{itemize}
  \item Abwehr als schonendstes Mittel (das Gelindeste der
  zur verlässlichen Abwehr verfügbaren Mittel)
  \item Zurechnungsprinzip (= richtet sich gegen den
  Angreifer)
  \item Die Annahme einer Notwehrsituation erfolgte bloß
  fahrlässig
\end{itemize}
% \end{minipage}
&
% \begin{minipage}{\linewidth}
\EE{Notwehrhandlung}:
\begin{itemize}
  \item Abwehr als schonendstes Mittel (das Gelindeste der
  zur verlässlichen Abwehr verfügbaren Mittel)
  \item Zurechnungsprinzip (= richtet sich gegen den Angreifer)
\end{itemize}
% \end{minipage}
\\\addlinespace\bottomrule
\end{tabularx}
\end{minipage}
\end{WidePar}
\end{table}




\begin{table}[H]\TableBaseModification\begin{WidePar}
\Captionof{table}{Rechtfertigender und
entschuldigender \EE{Notstand}\label{tab:notstand}}
\begin{minipage}{\linewidth}
\begin{tabularx}{\linewidth}{XX}\toprule\addlinespace
\noindent\II[Notstand!rechtfertigender]{Rechtfertigender Notstand} 
&
\noindent\II[Notstand!entschuldigender]{Entschuldigender
Notstand} \\\addlinespace\midrule\addlinespace
\begin{itemize}
    \item In keinem Gesetz aufgezeichnet; wird aber
    einheitlich angenommen
    \item Anzuwenden wenn die Gefahr auch tatsächlich
    (objektiv) vorliegt
\end{itemize}
&
\begin{itemize}
    \item Gerichtliches Strafrecht und Bürgerliches
    Recht (StGB, ABGB\ZFootnote{{\S}\,10 StGB und
    {\S}\,1306a ABGB})
    \item Anzuwenden wenn man die Gefahr bloß subjektiv
    annimmt
\end{itemize}
\\\addlinespace\midrule\addlinespace

\EE{Notstandssituation}: Gefahr ist:

\begin{itemize}
    \item  Kein menschlicher Angriff
    \item  Gegenwärtig oder unmittelbar drohend
    \item Rechtswidrig
    \item  Auf beliebiges Rechtsgut gerichtet
    \item  Droht bedeutender Nachteil
\end{itemize}

\EE{Notstandshandlung}:

\begin{itemize}
    \item  Abwehr als schonendstes Mittel
    \item  Zurechnungsprinzip (=~Richtet sich gegen die
    Gefahrenursache)
    \item  Geringes Risiko bei hoher Rettungschance 
\end{itemize}
&
\EE{Notstandssituation}: Gefahr ist:

\begin{itemize}
  \item  Kein menschlicher Angriff
  \item  Gegenwärtig oder unmittelbar drohend
  \item Rechtswidrig
  \item  Auf beliebiges Rechtsgut gerichtet
  \item  Droht bedeutender Nachteil
\end{itemize}

\EE{Notstandshandlung}:

\begin{itemize}
  \item  Abwehr als schonendstes Mittel
  \item  Zurechnungsprinzip (=~Richtet sich gegen die
  Gefahrenursache)
  \item  Drohender Nachteil darf geringfügig kleiner als das
  Ergebnis der Abwehrhandlung sein.
\end{itemize}
\\\addlinespace\bottomrule
\end{tabularx}
\end{minipage}\end{WidePar}
\end{table}



\begin{table}[H]\TableBaseModification\begin{WidePar}
\Captionof{table}{\EE{Anhalterecht}
Privater\label{tab:anhalterecht-privater}}

\begin{minipage}{\linewidth}
\begin{tabularx}{\linewidth}{X}\toprule\addlinespace
\noindent\II{Anhalterecht Privater}
\\\addlinespace\midrule\addlinespace
% \begin{minipage}{\linewidth}
\begin{itemize}
  \item Gerichtliches Strafrecht (StPO\ZFootnote{\S\,80 Abs. 2 StPO})
  \item Anzuwenden wenn ein Verdacht auf eine gerichtlich
  strafbare Handlung vorliegt
\end{itemize}
% \end{minipage}
\\\addlinespace\midrule\addlinespace
% \begin{minipage}{\linewidth}
\EE{Situation:}

\begin{itemize}
  \item Verdacht auf eine gerichtlich strafbare Handlung
\end{itemize}
% \end{minipage}

\\\addlinespace
% \begin{minipage}{\linewidth}
\EE{Handlung}:
\begin{itemize}
  \item Anhalten als schonendstes Mittel (in angemessener
  Weise)
  \item Zurechnungsprinzip (= Anhaltung des mutmaßlichen
  Täters)
  \item Unverzügliche Verständigung der Sicherheitspolizei
\end{itemize}
% \end{minipage}
\\\addlinespace\bottomrule
\end{tabularx}
\end{minipage}
\end{WidePar}
\end{table}

\Text{Fallbeispiele}
{
\begin{itemize}
  \item Fall 1: Randalierender Patient
  \dotfill\FullPrettyRef{topic:jus-langer-dienst-fall-1}
  \item Fall 2: Gefahr durch Hund
  \dotfill\FullPrettyRef{topic:jus-langer-dienst-fall-2}
\end{itemize}
}
{}
{}
{}
{}


\subsubsection{Unterbringungsrecht}

Das Unterbringungsrecht kommt bei unmittelbar eigen- oder
fremdgefärdenden psychiatrischen Patienten unter
bestimmten Voraussetzungen zur Anwendung.
Grundsätzlich gibt es beim Umgang mit solchen Patienten
mehrere Möglichkeiten:

% \Layout{22}{Möglichkeiten}{
\begin{enumerate}
  \item Freiwillige Behandlung (Einwilligung des Patienten,
  er behält alle Rechte)
  \item Unterbringung auf Verlangen (er
  verzichtet auf manche Rechte)
  \item Unterbringung ohne Verlangen (gegen oder ohne Willen
  des Patienten)
\end{enumerate}


Die \I{Unterbringung} wird durch das
\I{Unterbringungsgesetz} (\I{UbG}) geregelt.
In eine psychiatrische Abteilung darf nur untergebracht
werden, wer
\begin{itemize}
  \item an einer \EE{psychischen Krankheit} leidet und im
  Zusammenhang damit
\item \EE{sein Leben} oder seine \EE{Gesundheit} oder das Leben
  \EE{oder} die Gesundheit \EE{anderer} ernstlich und erheblich
  gefährdet und
  \item \EE{nicht in anderer Weise}, insbesondere außerhalb einer
  Anstalt, ausreichend ärztlich behandelt oder betreut
  werden kann.\ZFootnote{{\S}\,3}
\end{itemize}

\EE{Eine Unterbringung kann nur in einer psychiatrische
Abteilung erfolgen!}

% }{}
% {./images/shared/mann_mit_pistole.jpg}
% {Sie kommen in eine Wohnung und finden diese
% Situation vor \ldots Was machen Sie?}
% {Gabriel, AASS-CC-AT.at}



\Text{Unterbringung ohne Verlangen}
{%
Eine Person darf gegen oder ohne ihren Willen nur dann in
eine psychiatrische Abteilung gebracht werden, wenn ein im
\E{öffentlichen Sanitätsdienst stehender Arzt} oder ein
\I{Polizeiarzt} (\II{Amtsarzt}) sie untersucht und
bescheinigt, dass die Voraussetzungen der Unterbringung
vorliegen. In der Bescheinigung sind im einzelnen die Gründe
anzuführen, aus denen der Arzt die Voraussetzungen der
Unterbringung für gegeben erachtet.\ZFootnote{{\S}\,8 UbG}



Die \EE{Sicherheitspolizei} ist berechtigt und verpflichtet,
eine Person, bei der sie aus besonderen Gründen die
Voraussetzungen der Unterbringung für gegeben erachtet, zur
Untersuchung zum Amtsarzt zu bringen oder diesen
beizuziehen. Bescheinigt der Amtsarzt das Vorliegen der
Voraussetzungen für eine Unterbringung, so hat die
Sicherheitspolizei die betroffene Person in eine Anstalt zu
bringen oder dies zu veranlassen.\ZFootnote{{\S}\,9\,(1)
UbG}

\EE{Bei Gefahr im Verzug} kann die \EE{Sicherheitspolizei}
die betroffene Person
\EE{auch ohne Untersuchung und Bescheinigung in eine Anstalt
bringen}.\ZFootnote{{\S}\,9 Abs.\,2 UbG}

Der Arzt und die Sicherheitspolizei haben unter möglichster
Schonung der betroffenen Person vorzugehen und die notwendigen
Vorkehrungen zur Abwehr von Gefahren zu treffen. Sie haben,
soweit das möglich ist, mit psychiatrischen Einrichtungen
außerhalb einer Anstalt zusammenzuarbeiten und
\EE{erforderlichenfalls} den \EE{örtlichen Rettungsdienst
beizuziehen}.\ZFootnote{{\S}\,9 Abs.\,3 UbG} \EE{Für den
Transport ist demnach grundsätzlich die Polizei
zuständig!\ZFootnote{{\S}\,9~Abs.\,1)}}

Sollte der Rettungsdienst von sich aus einen
unterzubringenden Patienten vorfinden, so ist die
Sicherheitspolizei und der Amtsarzt zu verständigen.
\EE{Ein Notarzt ist i.\,d.\,R. \E{kein} Amtsarzt}, auch wenn
er von einer Gemeinde gestellt wird, und
hat keine Kompetenz, eine Unterbringung anzuordnen. 

\opt{REGVIE}{In Wien
ist der Amtsarzt über die Bundespolizei anzufordern.}


\Cave{Die Persönlichkeitsrechte psychisch Kranker, die in
eine Krankenanstalt aufgenommen werden, sind besonders zu
schützen. Die Menschenwürde psychisch Kranker ist unter
allen Umständen zu achten und zu
wahren.\footnote{{\S}\,1~Abs.\,1 UbG}}
}
{}
{}
{}
{}




\Text{Unterbringung mit Verlangen}
{
Eine Unterbringung \EE{kann auch freiwillig}, d.\,h. auf
Wunsch des Patienten erfolgen. Da dann auch der Transport
freiwillig stattfindet, sind dabei (beim Primäreinsatz)
keine besonderen Bestimmungen zu beachten.
}
{}
{}
{}
{}

\Text{Fallbeispiele}
{
\begin{itemize}
  \item Fall 1: Randalierender Patient
  \dotfill\FullPrettyRef{topic:jus-langer-dienst-fall-1}
\end{itemize}
}
{}
{}
{}
{}


%%%%%%%%%%%%%%%%%%%%%%%%%%%%%%%%%%%%%%%%%%%%%%%%%%%%%%%%%%%
%%%%%%%%%%%%%%%%%%%%%%%%%%%%%%%%%%%%%%%%%%%%%%%%%%%%%%%%%%%
%%%%%%%%%%%%%%%%%%%%%%%%%%%%%%%%%%%%%%%%%%%%%%%%%%%%%%%%%%%
%%%%%%%%%%%%%%%%%%%%%%%%%%%%%%%%%%%%%%%%%%%%%%%%%%%%%%%%%%%
%%%%%%%%%%%%%%%%%%%%%%%%%%%%%%%%%%%%%%%%%%%%%%%%%%%%%%%%%%%
\subsection{Rechte des Patienten}

\label{topic:patientenrechte}

\subsubsection{Gleichheitsgrundsatz}

\index{Gleichheitsgrundsatz}

Alle Patienten sind gleich zu behandeln, unabhängig
ihrer Religion, Hautfarbe, ethnische Zugehörigkeit,
geschlechtliche Orientierung oder ihres Geschlechts oder
sonstiger unsachlicher Unterscheidung.\ZFootnote{(Art.~7
B-VG, Art.~2 StGG, Art.~20 Z.~1 EGC)}

\subsubsection{Patientenrechte nach dem
Krankenanstalten- und Kuranstaltengesetz (KAKuG)}

Als zentrale Orientierungshilfe für Rechte eines Patienten
dient auch für das Rettungswesen die Aufzählung an
Patientenrechten des KAKuG\ZFootnote{{\S}\,5a KAKuG},
z.\,B.:

\begin{itemize}
  \item Informierung des Patienten über seine Rechte
  \item Einsichtnahme in die ihn betreffende
  medizinische Dokumentation
  \item Informierung des Patienten auf möglichst
  schonungsvolle und verständliche Art
  \item Wahrung der Intimsphäre
\end{itemize}


\subsubsection{Heilbehandlung}
\label{topic:heilbehandlung}

Der Begriff der \T{Heilbehandlung} wird aus allgemeinen
Rechtsgrundsätzen und systematischer Betrachtung der
gesamten Rechtsordnung gewonnen.

Es müssen folgende \EE{Voraussetzungen} vorliegend, damit
eine Handlung als Heilbehandlung gilt:
\begin{itemize}
    \item Medizinische Indikation (zur
    Diagnose oder Linderung von Erkrankungen und/oder
    Verletzungen notwendig)
    \item Durchführung ist lege artis (entspricht dem
    geltenden Stand der Wissenschaft)
    \item Aufklärung und Einwilligung des Patienten (sofern einsichts- und urteilsfähig; vgl.
    \FullPrettyRef{topic:jus-fall-aufklaerung})
\end{itemize}

Sobald diese Voraussetzungen vorliegen, wird eine
Heilbehandlung, auch wenn sie Einwirkung auf den Körper
des Patienten nimmt, als gerechtfertigt angesehen. 
Werden die Voraussetzungen der Heilbehandlung hingegen nicht
eingehalten, macht man sich u.\,U. straf- bzw.
haftbar\ZFootnote{\begin{inparaenum}
  \item \EE{Strafrecht}:
  \begin{inparaenum}
    \item \E{Körperverletzung} ({\S}{\S}\,83\,ff StGB)
    \item \E{Eigenmächtige  Heilbehandlung }
    ({\S}\,110 StGB,
    \FullPrettyRef{topic:eigenmächtige-heilbehandlung}).
  \end{inparaenum}
  \item \EE{Zivilrecht}:
  \begin{inparaenum}
    \item \E{Schadenersatzansprüche} ({\S}{\S}\,1293\,ff
    ABGB)
  \end{inparaenum}
\end{inparaenum}
}.

\Text{Eigenmächtige Heilbehandlung}
{\label{topic:eigenmächtige-heilbehandlung}
Von einer \I{eigenmächtigen Heilbehandlung} spricht man,
wenn der Patient in die Behandlung \E{nicht einwilligt} bzw.
\E{nicht aufgeklärt} wird. Dies ist nur bei \Speech{\enquote{Gefahr
in Verzug}} (\E{unmittelbar} drohende oder
gegenwärtige Gefahr, die ein Abwarten nicht zulässt bzw.
unmittelbares Handeln verlangt) zulässig. Sonst stellt
sie einen gerichtlichen \EE{Straftatbestand}
dar\ZFootnote{{\S}\,110 StGB}, da sie dem
Selbstbestimmungsrecht des Patienten zuwider läuft.


} {}{}{}{}

\Text{Fallbeispiele}
{
\begin{itemize}
  \item Fall 3: Blutzuckermessung als Heilbehandlung
  \dotfill\FullPrettyRef{topic:jus-langer-dienst-fall-3}
  \item Fall 4: Aufklärung und Einwilligung
  \dotfill\FullPrettyRef{topic:jus-langer-dienst-fall-4}
\end{itemize}
}
{}
{}
{}
{}

\subsubsection{Die Patientenverfügung}


Die Patientenverfügung ist die Willenserklärung einer
einsicht-, urteils- und äußerungsfähigen Person, die
erst dann Bedeutung erlangt, wenn der Patient seinen Willen
nicht mehr ausdrücken kann.
Dabei kann nur die Ablehnung einer bestimmten
Heilbehandlung Inhalt der Verfügung sein.

Ein auf jeden Fall \EE{unbeachtlicher Inhalt} wäre der
Ausschluss einer Grundversorgung oder eine medizinisch
nicht indizierte Behandlung. Genauso führen zur
Unbeachtlichkeit ein zu
allgemeine oder widersprüchliche Formulierungen, ein
rechtswidriger Inhalt oder erkennbare Willensmängel des
Patienten bei Abgabe der Verfügung (nicht ernst gemeint,
Irrtum, List, Drohung oder Täuschung)

\paragraph{Zwei Arten der Patientenverfügung}

Es gibt zwei Arten von Patientenverfügungen, eine
\E{verbindliche} und eine
bloß \E{beachtliche
Patientenverfügung}. 


\begin{enumerate}
  \item \T[Patientenverfügung!Verbindliche]{Verbindliche Patientenverfügung}:
  \begin{itemize}
    \item \EE{Schriftlich} vor Notar, Rechtsanwalt oder
    rechtskundigen Patientenvertreter
    \item Eigenhändig unterschrieben und datiert
    \item Aufgrund einer umfassenden Aufklärungen durch
    einen Arzt
    \item Inhalt muss eine \EE{konkrete Heilbehandlung}
    ablehnen ${\rightarrow}$ allgemeine oder
    widersprüchliche Formulierungen führen zur Unwirksamkeit
    \item Wirkung nur für \EE{5 Jahre}
    \item Eine Verfügung, die alle Elemente enthält ist
    \EE{absolut bindend} und \EE{darf nicht umgangen
    werden!}
  \end{itemize}
  \item \T[Patientenverfügung!Beachtliche]{Beachtliche Patientenverfügung}:
  \begin{itemize}
    \item
    \EE{grundsätzlich}: alle Verfügungen, die nicht
    die obigen Voraussetzungen erfüllen
    \item
    Umfang der Beachtlichkeit: \EE{bloß relative
    Beachtlichkeit}:
    \begin{itemize}
      \item
      Verfügung soll \EE{in Behandlungsentscheidungen
      berücksichtigt} werden, aber
      \item
      das Wohl des Patienten bleibt oberstes Gebot
      \item
      daher kann man sich bei besonderer Begründung über die Verfügung
      \EE{hinwegsetzen}
    \end{itemize}
  \end{itemize}
\end{enumerate}

\Text{Suche nach Patientenverfügung}
{
Bei der Notfallversorgung darf die Suche nach einer
Verfügung auf keinen Fall die Versorgung des Patienten
beeinträchtigen. Solange nicht \EE{zweifelsfrei} feststeht,
dass eine Patientenverfügung verbindlich ist, muss der
Patient entsprechend behandelt werden.

Sollte das Personal nicht ausreichen, muss eine Suche
unterbleiben.\ZFootnote{{\S}\,12 PatVG} Auch das Lesen der
Verfügung bzw. die Abklärung, ob die Verfügung verbindlich
oder beachtlich ist, darf die Versorgung nicht
beeinträchtigen.

Das bedeutet für den Rettungsdienst: 
\begin{enumerate}
  \item Bei Notfallversorgungen darf nicht nach einer
  Verfügung gesucht werden, wenn die Suche das Leben oder
  die Gesundheit des Patienten ernstlich gefährden würde.
  \item \EE{Aber:} Wenn \E{zweifelsfrei} eine verbindliche
  Patientenverfügung vorliegt (z.\,B. weil eine betreuende
  Pflegekraft diese in der Zwischenzeit gefunden und die
  Verbindlichkeit bestätigt hat), muss dieser Folge
  geleistet werden.
\end{enumerate}

Liegt eine beachtliche Patientenverfügung vor, so ist im
Einzelfall zu entscheiden: Das Wohl des Patienten bleibt
oberstes Gebot und die Verfügung soll in die
Entscheidung einfließen. I.\,d.\,R. wird man bis zur Klärung
der genauen Umstände eine Reanimation beginnen bzw.
fortführen müssen.

}
{}
{}
{}
{}



\Text{Fallbeispiele}
{
\begin{itemize}
  \item Fall 5: Patientenverfügung
  \dotfill\FullPrettyRef{topic:jus-langer-dienst-fall-5}
\end{itemize}
}
{}
{}
{}
{}

%%%%%%%%%%%%%%%%%%%%%%%%%%%%%%%%%%%%%%%%%%%%%%%%%%%%%%%%%%%
%%%%%%%%%%%%%%%%%%%%%%%%%%%%%%%%%%%%%%%%%%%%%%%%%%%%%%%%%%%
%%%%%%%%%%%%%%%%%%%%%%%%%%%%%%%%%%%%%%%%%%%%%%%%%%%%%%%%%%%
%%%%%%%%%%%%%%%%%%%%%%%%%%%%%%%%%%%%%%%%%%%%%%%%%%%%%%%%%%%
%%%%%%%%%%%%%%%%%%%%%%%%%%%%%%%%%%%%%%%%%%%%%%%%%%%%%%%%%%%
\subsection{Berufs- und Tätigkeitsrecht}

%%%%%%%%%%%%%%%%%%%%%%%%%%%%%%%%%%%%%%%%%%%%%%%%%%%%%%%%%%%
%%%%%%%%%%%%%%%%%%%%%%%%%%%%%%%%%%%%%%%%%%%%%%%%%%%%%%%%%%%
%%%%%%%%%%%%%%%%%%%%%%%%%%%%%%%%%%%%%%%%%%%%%%%%%%%%%%%%%%%
\subsubsection{Die Aufgaben des Sanitäters}

Der Sanitätsdienst umfasst den Tätigkeitsbereich des
Rettungssanitäters und des Notfallsanitäters entsprechend
die eigenverantwortliche Anwendung von Maßnahmen der 
qualifizierten Ersten Hilfe, Sanitätshilfe und 
Rettungstechnik, einschließlich diagnostischer und
therapeutischer Verrichtungen.\ZFootnote{\S\,8~SanG}

Sanitäter haben ihre Tätigkeit ohne Ansehen der
Person gewissenhaft auszuüben. Sie haben das Wohl der
Patienten und der betreuten Personen nach Maßgabe der
fachlichen und wissenschaftlichen Erkenntnisse und
Erfahrungen zu wahren. Nötigenfalls ist ein Notarzt oder,
wenn ein solcher nicht zur Verfügung steht, ein sonstiger
zur selbständigen Berufsausübung berechtigter Arzt
anzufordern. Sanitäter haben sich tätigkeitsrelevant
fortzubilden.\ZFootnote{\S\,4~SanG} 


\Text{Fallbeispiele}
{
\begin{itemize}
  \item Fall 6: Unterlassen der Notarzt-Nachforderung
  \dotfill\FullPrettyRef{topic:jus-langer-dienst-fall-6}
  \item Fall 8: Rechtswidrige Anordnung durch Notarzt
  \dotfill\FullPrettyRef{topic:jus-langer-dienst-fall-8}
\end{itemize}
}
{}
{}
{}
{}

%%%%%%%%%%%%%%%%%%%%%%%%%%%%%%%%%%%%%%%%%%%%%%%%%%%%%%%%%%%
%%%%%%%%%%%%%%%%%%%%%%%%%%%%%%%%%%%%%%%%%%%%%%%%%%%%%%%%%%%
%%%%%%%%%%%%%%%%%%%%%%%%%%%%%%%%%%%%%%%%%%%%%%%%%%%%%%%%%%%
\subsubsection{Pflichten des Sanitäters}
\label{topic:sanitaeter-pflichten}

\begin{itemize}
  \item Behandlungs- bzw. Hilfeleistungspflicht  
  \item Dokumentationspflicht
  \item Verschwiegenheitspflicht
  \item Auskunftspflicht
  \item Fortbildungspflicht
  \item Rezertifizierungspflicht
\end{itemize}

\Text
{Behandlungs- bzw. Hilfeleistungspflicht}
{
\subparagraph{Ersthelferische Versorgungspflicht für
Durchschnittsbürger} Jeder Bürger (und nicht mehr
tätigkeitsberechtigte Sanitäter) ist zur Setzung von
Hilfsleistungen verpflichtet, die seinem Wissens- und
Kenntnisstand entsprechen. Daher ist jeder mit einem
Führerschein-Erste-Hilfe-Kurs zur Setzung von zumutbaren
Erste-Hilfe Maßnahmen verpflichtet.\ZFootnote{{\S}\,95 StGB}

\subparagraph{Behandlungspflicht für Sanitäter}
Jeder tätigkeits- oder berufsberechtigte Sanitäter,
unabhängig von seinem Kompetenzstand, ist zur
selbstständigen und eigenverantwortlichen Versorgung einer
hilfsbedürftigen Person verpflichtet. Dabei muss er, auch
außerhalb des Dienstes, zumindest qualifizierte Erste Hilfe
leisten und wenn nötig einen Arzt (jeder mit jus
practicandi) nachfordern.\ZFootnote{{\S}\,4 i.\,V.\,m.
{\S}\,8 i.\,V.\,m. {\S}\,9 Abs. 1 Z 1 SanG}

\subparagraph{Ärztliche Behandlungspflicht}
Ärzte, unabhängig ihrer Fachausbildung, sind dazu
verpflichtet, hilfsbedürftigen Patienten ärztliche
Grundversorgung zukommen zu lassen und -- wenn nötig --
weitere Einsatzkräfte nachzufordern oder den Patienten in
weitere medizinische Behandlung zu
schicken.\ZFootnote{{\S}\,48 ÄrzteG}

\subparagraph{Behandlungspflicht für Krankenanstalten}
Jede Krankenanstalt, ob öffentlich oder privat, ist
verpflichtet einen behandlungspflichtigen Patienten 
aufzunehmen und zumindest ärztliche Grundbehandlung zukommen
zu lassen.\ZFootnote{{\S}{\S}\,22 und 40 KAKuG}

Das in Wien geltende \E{Bettenkontingentsystem}, mit seinen
Zuweisungen zu bestimmten Krankenanstalten, ist trotzdem
zulässig, da sich der Patient während des Transportes schon
in einem präklinischen Umfeld (= Einsatzwagen mit Sanitätern
und Medizinprodukten) befindet. Sollte sich aber der Zustand
des Patienten während der Fahrt verschlechtern, darf das
nächstgelegene Spital angefahren werden.
Erst wenn der Patient wieder transportfähig ist, darf er mit
geeigneten Krankentransportmitteln in eine andere
Krankenanstalt weitergeschickt werden.

\subparagraph{Unzumutbarkeit der Behandlung}
Für alle Berufe oder Durchschnittsbürger ist die Behandlung
dann nicht zumutbar, wenn man \ldots
\begin{itemize}
  \item das eigene Leben oder die Gesundheit gefährden würde, oder
  \item höherwertige Interessen verletzen
  müsste.\ZFootnote{{\S}\,95 Abs. 2 StGB}
\end{itemize}

}
{%
\begin{itemize}
  \item Ersthelferische Versorgungspflicht für
  Durchschnittsbürger
  \item Ärztliche Behandlungspflicht
  \item Behandlungspflicht für Sanitäter 
  \item Behandlungspflicht für Krankenanstalten
  \item Nicht bei Unzumutbarkeit
\end{itemize}}
{}
{}
{}

\Text
{Dokumentationspflicht}
{\label{topic:dokumentationspflicht}%
\label{topic:dokumentation-rechtlich}%
Schon nach dem Sanitätergesetz besteht für jeden Sanitäter
die Pflicht, die von ihm gesetzten
sanitätsdienstlichen Maßnahmen zu
dokumentieren.\ZFootnote{{\S}\,5 SanG} Darüber hinaus müssen
analog zu anderen Gesundheitsberufen auch alle anderen
medizinisch relevanten Sachverhalte schriftlich und
nachvollziehbar fixiert werden.

Hintergrund dieser Dokumentationspflicht ist der
\EE{Heilbehandlungsvertrag}.
Im österreichischen Recht hat grundsätzlich der Kläger den
Schaden, die Rechtswidrigkeit und das Verschulden des
Beklagten zu beweisen. Besteht aber zwischen Kläger und
Beklagten ein Vertrag, so wird gem. \Z{{\S}\,1298} ABGB das
Verschulden des Beklagten angenommen bzw. der Beklagte muss
beweisen, dass ihn kein Verschulden trifft
(\T{Beweislastumkehr}).
Um nun für den Beklagten die Möglichkeit zu schaffen, sein
Unverschulden zu beweisen, wurde die
\I{Dokumentationspflicht} eingeführt, mit der bei
ausreichender Dokumentation wieder der Kläger das
Verschulden des Beklagten zu beweisen hat.


Inhalt der Dokumentation sollte nicht nur die 
gesetzten Maßnahmen sein, sondern auch 
\begin{compactitem}
  \item situationsrelevante
  Umstände (beim Patient; am BO; während des Transportes;
  {\dots}),
  \item  u.\,U. relevante Inhalte einer Aufklärung,
  \item Einwilligung bzw. Ablehnung durch den Patienten
  (Revers) sowie
  \item sonstige für den Sanitäter entscheidend wirkende
  Umstände.
\end{compactitem}




% \FazitLike{Grundsatz: Alles darf dokumentiert werden, es darf aber nicht Nichts
% dokumentiert werden.}

Es muss \E{gründlich}, \E{ausführlich} und
\E{nachvollziehbar} dokumentiert werden.
Die Dokumentation soll vorrangig am Einsatzprotokoll selbst
oder, wenn nötig, auf einem beigelegten
Blatt Papier erfolgen.


% Des Weiteren ist eine zusätzliche eigene Dokumentation für jeden
% Sanitäter empfehlenswert.



}
{}
{}
{}
{}

\Text
{Verschwiegenheitspflicht}
{%
\label{topic:verschwiegenheitspflicht}
\T{Geheimnisse}, die man aufgrund seiner
Tätigkeit als Arzt oder Sanitäter anvertraut bekommen
hat, dürfen nicht an unbefugte Dritte weitergegeben
werden.
{Geheimnisse} sind alle Informationen, die
einen Bezug zum Patienten haben oder liefern (z.\,B.
Name des Patienten, Versicherungsnummer, Wohnadresse,
{\dots}).

Geheimnisse \E{dürfen} unter Umständen
mitgeteilt werden (\T{Auskunftsrecht}), z.\,B.
 wenn die Anonymität des Patienten gewahrt bleibt. 
Daher darf man mit Kollegen oder sonstigen Personen
über die erlebten Geschehnisse reden, solange aus dem
Gesagten nicht erkennbar ist, um welche Person es sich
bei dem Patienten handelt (\E{Anonymisierung}).
Ebenso dürfen Geheimnisse weitergegeben werden, wenn der
Patient das Fachpersonal von der \E{Schweigepflicht
entbunden} hat.
}
{}
{}
{}
{}

\Text
{Informationspflichten}
{%

\begin{itemize}
  \item \T{Auskunftspflicht}: \label{topic:auskunftspflicht}
  Geheimnisse \E{müssen}
  unter bestimmten Umständen weitergegeben werden. Dies
  geschieht für Sanitäter im Rahmen der sog.
  \IX{Auskunftspflicht}\ZFootnote{{\S}\,7 SanG} in folgenden
  Fällen:
  \begin{itemize}
    \item   Wenn es zur Behandlung des Patienten dient
    \item   Wenn es gesetzlichen vorgeschrieben ist (z.\,B. im
    Gerichtsverfahren als Zeuge)
    \item Zum Schutz von höherwertigen Interessen der
    öffentlichen Gesundheitspflege oder Rechtspflege
    (${\rightarrow}$ zur Meldepflicht
    siehe unten)
  \end{itemize}
  Diese Auskunftspflicht besteht gegenüber
  \begin{itemize}
    \item   dem Patienten
    \item   gesetzlichen Vertreter des Patienten
    \item   anderen Personen der Gesundheitsberufe, die den
    Patienten betreuen, behandeln oder pflegen bzw. an denen
    der Patient im Zuge des Transportes übergeben wird \item
    Sozialversicherung und Krankenanstalt \item vom Patient als
    auskunftsberechtigt benannte Personen
  \end{itemize}
  \item \T{Meldepflicht}: Bestimmte leicht
  übertragbaren Krankheiten, bei denen die Gefahr von
  schweren Gesundheitsschäden und/oder Todesfolgen besteht,
  sind einer Behörde zu melden (\EE{Meldepflichtige
  Krankheiten}). Grundlage dieser Meldepflicht ist stets ein
  Gesetz oder eine Verordnung, wobei als zentrale rechtliche
  Bestimmungen folgende Normen gelten:\Z{
  \begin{itemize}
    \item \I{Epidemiegesetz}: Bei bloßen Verdacht oder
    sicherer Diagnose einer im Gesetz genannten Erkrankung.
    \item \I{Tuberkulosegesetz}: bei bloßen Verdacht auf
    eine Erkrankung die mit Sicherheit oder
    Wahrscheinlichkeit durch Tuberkulosebakterien verursacht
    wurde
   \item \I{Geschlechtskrankheitengesetz}: nur bei
   Befürchtung, dass sich die Krankheit weiterverbreiten
   würde oder sich der Patient einer Heilbehandlung
   entziehen würde.
   \item \I{AIDS-Gesetz}: \E{Anonymisierte} Meldepflicht
   erst ab manifester Erkrankung mit AIDS (=~bei sichtbaren
   Krankheitszeichen; nicht schon bei HIV-Infektion) oder
   Todesfall mit erkennbaren AIDS typischen
   Krankheiten\footnote{s.
   \prettyref{topic:aids}}.
   Es gibt \E{keine} Untersuchungs- oder Behandlungspflicht,
   daher darf kein behördlicher Zwang ausgeübt werden \item
   nach \EE{Verordnungen} des Gesundheitsministerium:
   Da sich ständig neue übertragbare Krankheiten\footnote{in
   letzter Zeit z.\,B.
   \E{SARS}, \E{Vogelgrippe}, \E{Schweinegrippe}
   (\prettyref{topic:grippe})} entwickeln, gibt das
   Gesundheitsministerium in unregelmäßigen  Abständen
   Verordnungen heraus, die eine Meldepflicht beinhalten.
  \end{itemize}
  
  Nicht jeder Durchschnittsbürger ist zur Meldung
  verpflichtet, sondern nur besondere Personengruppen
  worunter aber alle Gesundheitsberufe fallen.
  
  Die Meldung hat bei AIDS an das Bundesministerium für
  Gesundheit zu ergehen, ansonsten an die zuständige
  Bezirksverwaltungsbehörde (in Wien an die MA\,15).
  }
  \item \T{Anzeigepflicht}:\label{topic:anzeigepflicht}
  Nur bestimmte Personengruppen sind dazu verpflichtet bei Kenntnis über
  strafbare Handlungen dies der Polizei oder Staatsanwaltschaft per
  Anzeige zu melden.
  
  \begin{itemize}
    \item \EE{\E{Ärztliche}
    Anzeigepflicht}\ZFootnote{{\S}\,54 Abs.
    4 und 5 ÄrzteG}:
    Nur bei Vorliegen folgender Delikte (Verdacht reicht aus) und
    unabhängig vom Willen des Betroffenen:
    \begin{itemize}
      \item Tötung bzw. schwere Körperverletzung
      \item Misshandlung, Quälen, Vernachlässigung oder
      sexueller Missbrauch
      von:
      \begin{itemize}
        \item Minderjährigen oder
        \item Volljährigen, die nicht alleine in der Lage
        sind ihre Interessen
        wahrzunehmen
      \end{itemize}
    \end{itemize}
    \item Polizeiliche Ermittlungspflicht\ZFootnote{{\S}\,99 StPO}
    für alle Delikte
  \item \T[Anzeigerecht]{Anzeige\underline{recht}}:
  \label{topic:anzeigerecht}
  Jeder Durchschnittsbürger und auch sonstige Gesundheitsberufe sind
  \E{nicht verpflichtet} zur Anzeige von Strafdelikten. Sie
  haben aber das
  \E{Recht}, eine strafbare Handlung anzuzeigen.\ZFootnote{{\S}\,80 Abs. 1 StPO}
  
  Trotz der fehlenden Verpflichtung für den Sanitäter sollte
  bei jedem Einsatz, bei dem der Verdacht einer strafbaren
  Handlung vorliegt, zumindest die Polizei nachgefordert
  werden. Die Nachalarmierung kann man in weiter Analogie
  mit der aus \Z{{\S}\,4} SanG entstammenden Pflicht für den
  Sanitäter
  \Speech{\enquote{das Wohl der Patienten [{\dots}] zu
  wahren}} begründen.
  \end{itemize}
\end{itemize}





}
{}
{}
{}
{}

\TextSlide
{Fortbildungspflicht}
{\label{topic:fortbildungspflicht}
Die Fortbildungspflicht für Sanitäter ist unter zwei
verschiedenen Gesichtspunkten zu sehen: Einerseits muss er
sich fachbezogen nach dem Sanitätergesetz, andererseits
gerätebezogen nach dem Medizinproduktegesetz (\Abk{MPG})
fortbilden.

\begin{itemize}
  \item \EE{Fortbildungspflicht nach
  SanG\ZFootnote{{\S}\,50 SanG}:}
  Jeder Sanitäter ist dazu verpflichtet, sich \EE{innerhalb
  von 2 Jahren} ab seinem Stichtag im Umfang von \EE{16
  Stunden fortzubilden}. Als \II{Stichtag} wird der
  Monatserste nach bestandener Rettungssanitäterausbildung
  gewertet, wobei sich der Stichtag auch bei Erreichen von
  höheren Kompetenzniveaus nicht ändert. Zu welchem Zeitpunkt
  die 16 Fortbildungsstunden innerhalb des Rahmens von 2
  Jahren besucht werden, ist dem Sanitäter überlassen.
  
  Man kann also entweder regelmäßig Fortbildungen besuchen
  oder auf einmal die Stundenzahl ansammeln. Weiters muss ein
  Sanitäter keine bestimmten Themen als Fortbildungen
  besuchen, sondern kann nach seinem Interesse vorgehen.
  
  \item \EE{Fortbildungspflicht nach
  MPG\ZFootnote{{\S}\,83 MPG}:}
  Die im Rettungswesen verwendeten Materialien gelten,
  abgesehen von Arzneimitteln, als Medizinprodukte und dürfen
  daher iSd. MPG nur von eingeschulten Personal angewendet
  werden. Diese Einschulungen müssen nicht nur alle
  Medizinprodukte der Rettungsorganisation beinhalten, sondern
  auch von allen Sanitätern besucht werden. Daher kann ein
  Sanitäter in Bezug auf Stundenanzahl oder Fortbildungsthemen
  nicht frei wählen. Die Fortbildungsstunden für das MPG
  können aber als Fortbildungsstunden iSd. SanG gewertet
  werden. Zusätzlich ist zu sagen, dass eine Einschulung der
  meisten Medizinprodukte schon durch die Standardausbildung
  zum Rettungssanitäter abgedeckt ist. Für die anderen,
  insbesondere höhere technische Produkte, muss jede
  Organisation separate Fortbildungen anbieten.
\end{itemize}



}
{\begin{itemize}
  \item Fortbildungspflicht SanG: 16\Hour* in 2 Jahren
  \item Fortbildungspflicht MPG: Medizinprodukte dürfen nur
  von eingeschultem Personal verwendet werden
\end{itemize}
}
{}
{}
{}

\TextSlide
{Rezertifizierungspflicht}
{\label{topic:rezertifizierungspflicht}%
{\ldots} besteht nach dem SanG\ZFootnote{{\S}\,51
SanG}: 
Innerhalb des 2 Jahre Rahmens der Fortbildungspflicht muss
sich ein Sanitäter, unabhängig von seinem Ausbildungsstand,
prüfen lassen in der Herz-Lungen-Wiederbelebung mit
halbautomatischen Defibrillator.
Der Rezertifizierungszeitpunkt innerhalb des 2 Jahre Rahmens
kann vom Sanitäter frei gewählt werden.

\noindent\EE{Sonderfall NKI}: Ein Notfallsanitäter mit
Kompetenzen der Intubation muss sich zusätzlich innerhalb seines 2
Jahre Rahmens einmal über die Intubation prüfen lassen.
}
{
\begin{itemize}
  \item Rezertifizierungspflicht SanG:
  Herz-Lungen-Wiederbelebung mit halbautomatischen
  Defibrillator, alle 2 Jahre
  
  NKI: Rezertifizierung der Kompetenz alle 2 Jahre
\end{itemize}
}
{}
{}
{}


\Text{Fallbeispiele}
{
\begin{itemize}
  \item Fall 9: Dokumentationspflicht
  \dotfill\FullPrettyRef{topic:jus-langer-dienst-fall-9}
  \item Fall 10: Verschwiegenheits- und Auskunftspflicht
  \dotfill\FullPrettyRef{topic:jus-langer-dienst-fall-10}
  \item Fall 11: Fortbildungs- und Rezertifizierungspflicht
  \dotfill\FullPrettyRef{topic:jus-langer-dienst-fall-11}
  \item Fall 14: Behandlungspflicht
  \dotfill\FullPrettyRef{topic:jus-langer-dienst-fall-14}
\end{itemize}
}
{}
{}
{}
{}





\subsubsection{Kompetenzstufen}
% Freigabe Tabelle KOLC 2010-08-25
\Layout{57}{}{}{}{
\begin{tabulary}{\linewidth}{LL}\toprule\addlinespace
\noindent\TabHeading{Kompetenzstufe}
&
\noindent\TabHeading{Beschreibung}
\\\addlinespace\toprule\addlinespace
\TabSub{RS}

\I{Rettungssanitäter}
&
\begin{minipage}[t]{\linewidth}
\begin{itemize}
  \item Selbstständige und eigenverantwortliche Versorgung
  und Betreuung
    des Patienten
    \item Übernahme und Übergabe des Patienten im Rahmen
    eines Transportes
    \item Hilfestellung bei Akutsituationen; Verabreichung von \ce{O2}
    \item Durchführung von lebensrettenden Sofortmaßnahmen
    \item Durchführung von Sondertransporten
    \item Blutzuckermessung
\end{itemize}
\end{minipage}


\\\addlinespace\toprule\addlinespace
\TabSub{NFS}

\I{Notfallsanitäter}
&
\begin{minipage}[t]{\linewidth}
\begin{itemize}
  \item Tätigkeiten des RS
  \item Unterstützung des Arztes; einschließlich der
  Betreuung und des sanitätsdienstlichen Transports von \E{Notfallpatienten}
  \item Verabreichung von Medikamenten der der
  \EE{Arzneimittelliste 1}:
    Wird vom leitenden Arzt der Rettungsorganisation
    festgelegt. Die Verabreichung kann durchgeführt werden,
    ohne dass ein Notarzt informiert werden muss, oder es
    ein solcher anordnet.
    \item Eigenverantwortliche Betreuung von berufsspezifischen
    Medizinprodukten und Arzneimitteln
    \item Mitarbeit in Forschung
\end{itemize}
\end{minipage}

\\\addlinespace\midrule\addlinespace
\TabSub{NKA}

Notfallsanitäter mit \E{allgemeiner Notkompetenz Arzneimittellehre}
&
\begin{minipage}[t]{\linewidth}
\begin{itemize}
  \item Zusätzlich zu den oben genannten Kompetenzen darf der NKA jene
  Medikamente der Arzneimittelliste 2 verabreichen, die ohne venösen
  Zugang verabreicht werden können.
  \item Die \EE{Arzneimittelliste 2} wird vom leitenden Arzt
  der Rettungsorganisation festgelegt. Die Medikamente dürfen
  auf Anweisung eines anwesenden Arztes (jeder mit jus
  practicandi, nicht unbedingt Notarzt) \E{oder} nach
  Verständigung eines Notarztes verabreicht werden.
\end{itemize}
\end{minipage}
 \\\addlinespace\midrule\addlinespace
\TabSub{NKV}

Notfallsanitäter mit \E{allgemeiner Notkompetenz Venenzugang und Infusion}
& 
\begin{minipage}[t]{\linewidth}
\begin{itemize}
  \item Zusätzlich zu den oben genannten
  Kompetenzen darf der NKV jene Medikamente der Arzneimittelliste 2
  verabreichen, die über venösen Zugang verabreicht werden können.
\end{itemize}
\end{minipage}

\\\addlinespace\toprule\addlinespace
\TabSub{NKI}

Notfallsanitäter mit \emph{besonderer Notkompetenz Beatmung und Intubation} 
&
\begin{minipage}[t]{\linewidth}
\begin{itemize}
  \item Durchführung einer endotrachealen Intubation samt
  Beatmung, aber ohne
  Prämedikation.
  \item Die NKI muss alle 2 Jahre rezertifiziert werden.
  \item Durchführung der Intubation auf Anweisung eines
  anwesenden Arztes (jeder mit jus
  practicandi, nicht unbedingt Notarzt) \E{oder} nach
  Verständigung eines Notarztes.
\end{itemize}
\end{minipage}




\\\addlinespace\bottomrule
\end{tabulary}
}
{Kompetenzstufen von Sanitätern nach dem SanG ({\S}{\S}\,\VonBis{9}{12})\label{tab:sang-kompetenzstufen}}
{} 

\Text{Fallbeispiele}
{
\begin{itemize}
  \item Fall 7: Sanitäterkompetenzen
  \dotfill\FullPrettyRef{topic:jus-langer-dienst-fall-7}
\end{itemize}
}
{}
{}
{}
{}




%%%%%%%%%%%%%%%%%%%%%%%%%%%%%%%%%%%%%%%%%%%%%%%%%%%%%%%%%%%
%%%%%%%%%%%%%%%%%%%%%%%%%%%%%%%%%%%%%%%%%%%%%%%%%%%%%%%%%%%
%%%%%%%%%%%%%%%%%%%%%%%%%%%%%%%%%%%%%%%%%%%%%%%%%%%%%%%%%%%
%%%%%%%%%%%%%%%%%%%%%%%%%%%%%%%%%%%%%%%%%%%%%%%%%%%%%%%%%%%
%%%%%%%%%%%%%%%%%%%%%%%%%%%%%%%%%%%%%%%%%%%%%%%%%%%%%%%%%%%
\subsection{Im Straßenverkehr}

\label{topic:einsatzfahrt}
\label{topic:einsatzfahrtt-missbrauch}
Im Straßenverkehr genießen Rettungsfahrzeuge eine besondere
Stellung, da sie Einsatzfahrten als Einsatzfahrzeuge
durchführen können.
Einsatzfahrzeuge sind Fahrzeuge, welche mit Blaulicht und
Folgetonhorn ausgestattet sind für die Dauer der Verwendung
dieser Signale.\footnote{Die StVO verlangt zwar nur, das
eines der beiden Signale benutzt wird, dennoch trifft den
Fahrer eines Einsatzfahrzeiges eine erhöhte
Sorgfaltspflicht. Dadurch ist die Verwendung beider Signale
gleichzeitig, zumindest in Gefahrensituationen und bei
Inanspruchnahme eines Vorrangs (rote Ampel, STOP-Tafel!),
geboten.} Eine Einsatzfahrt
darf nur bei Gefahr im Verzug vorgenommen werden, das
Blaulicht darf ferner auch zur Absicherung am
Ort einer Hilfeleistung eingesetzt
werden.\ZFootnote{{\S}\,26~StVO} 
Wer Blaulicht oder Folgetonhorn \E{missbräuchlich}
verwendet, macht sich gerichtlich strafbar nach dem
\E{Bundesgesetz gegen den Missbrauch von Notzeichen}. 

\subparagraph{Verhalten bei Blaulichtfahrt}

\begin{itemize}
    \item Es dürfen weder Personen noch Sachen gefährdet oder beschädigt
    werden
    \item Verkehrsgebote und -beschränkungen gelten
    grundsätzlich nicht -- Naturgesetze
    schon!\footnote{Denke an Höhen-, Breiten- und Gewichtsbeschränkungen!}
    \item Geschwindigkeitshöchstgrenzen gelten nicht.
    \item Ampeln, wenn sie nicht auf grün gschaltet sind,
    gelten als STOP-Tafeln.
    \item Absoluter Vorrang vor Nicht-Einsatzfahrzeugen.
    \item Es darf kein Nicht-Einsatzfahrzeug dem Rettungswagen unmittelbar
    nachfahren.
    \item \Ca{Der Vertrauensgrundsatz ist
    ungültig}: \EE{Man darf bei Blaulichtfahrten nicht mehr
    darauf vertrauen, dass sich die anderen
    Verkehrsteilnehmer richtig verhalten!} Daher ist
    defensives und voraussehendes Fahren verpflichtend, es
    gilt eine erhöhte Sorgfaltspflicht.
\end{itemize}

\Fazit{%
Nur wer noch größere Sorgfalt und Vorsicht wahrt als bei
normaler Fahrt, verhält sich richtig!
\\
{Das beinhaltet mitunter eine \EE{geringere}
Fahr\-ge\-schwin\-dig\-keit als die höchst zulässige!} }

\subparagraph{Exkurs: Anschnallpflicht}

Der Fahrer von jedem Fahrzeuges ist dazu verpflichtet zu
achten, dass er und seine Passagiere angeschnallt
sind.\ZFootnote{{\S}\,106~KFG} Es steht aber jedem Passagier
frei zu sagen, dass er unangeschnallt mitfahren möchte,
wodurch eine Schuld beim Fahrer entfällt. Für Kranken- und
Rettungstransport besteht zusätzlich gem.
Kraftfahrzeuggesetz (KFG) die Ausnahme, dass ein Passagier
(Personal wie Patient) unangeschnallt mitfahren darf, wenn
es ansonsten mit dem Zweck der Fahrt unvereinbar
ist.\ZFootnote{{\S}\,106 Abs.\,3~Z\,3 KFG} Trotzdem sollte
der Patient aufgefordert werden sich anzuschnallen bzw.
sollte vom Sanitäter angeschnallt werden.
Verweigert der Patient, muss diese Verweigerung dokumentiert
werden (am besten mittels Revers, den der Patient
unterschreibt).

Wird man nun im Zuge eines (sogar unverschuldeten)
Verkehrsunfalls verletzt und war man dabei nicht
angeschnallt, verliert man einen Teil des
Schandenersatzanspruchs, weil das Fahren ohne angelegten
Gurt als Mitverschulden gewertet wird.

\subparagraph{Vorrang der Einsatzfahrzeuge untereinander:}

\begin{enumerate}
    \item Rettungsfahrzeuge 
    \item Feuerwehrfahrzeuge
    \item Polizeifahrzeuge
    \item sonstige Einsatzfahrzeuge 
\end{enumerate}


\subparagraph{Ohne Blau geht's auch}

Fahrzeuge, welche über Blaulicht und Folgetonhorn verfügen,
dürfen auch außerhalb von Einsatzfahrten Fahrspuren und
Straßen für Omnibusse befahren. Weiters dürfen sie auch das
allgemeine Fahrverbot missachten und gegen die
vorgeschriebene Fahrtrichtung von Einbahnstraßen
fahren, wenn bereits Ausnahmen für
andere Fahrzeuge oder Fuhrwerke
bestehen.\ZFootnote{{\S}\,26a StVO}



\Text{Fallbeispiele}
{
\begin{itemize}
  \item Fall 17: Blaulichtverwendung 
  \dotfill\FullPrettyRef{topic:jus-langer-dienst-fall-17}
\end{itemize}
}
{}
{}
{}
{}


\ChapterEnd
